%%%%
\teteSndMole

%%%% titre
\vspace*{-36pt}
\titreActivite{En quête de stabilité : les ions}


%%%% Objectifs
\vspace*{-8pt}
\begin{objectifs}
  \item Comprendre la règle du duet et de l'octet.
  \item Comprendre comment se forment les ions à partir des atomes neutres.
\end{objectifs}

\begin{contexte}
  Dans la nature la plupart des atomes vont spontanément perdre ou gagner des électrons pour former des ions ou des molécules.
  Les seuls éléments qui ne réagissent pas et qu'on trouve sous forme de gaz monoatomique, sont  les gaz nobles de la 18$^\text{ème}$ colonne du tableau périodique (\chemfig{He}, \chemfig{Ne}, \chemfig{Ar}, \chemfig{Kr}, etc.).
  C'est parce qu'ils ont une grande stabilité, on dit qu'ils ont une grande inertie chimique.
  
  \problematique{
    Comment expliquer la formation d'ions monoatomique et la charge qu'ils portent à partir de la configuration électronique des gaz nobles ?
  }
\end{contexte}


%%%% question
\numeroQuestion Compléter le tableau suivant

\begin{center}
  
  \begin{tableau}{|c | c | c | c |}
     Gaz noble &
     Numéro atomique & Nombre d'électrons &
     Configuration électronique \\
     Hélium \chemfig{He} & $Z = 2$  & \correction{2}  & \correction{$1s^2$} \\
     Néon \chemfig{Ne}   & $Z = 10$ & \correction{10} & \correction{$1s^2 2s^2 2p^6$} \\
     Argon \chemfig{Ar}  & $Z = 18$ & \correction{18} & \correction{$1s^2 2s^2 2p^6 3s^2 3p^6$} \\
  \end{tableau}
\end{center}

\question{
  Comment est la couche externe pour ces trois gaz nobles ?
}{
  Leur couche externe (1, 2 ou 3) est pleine.
}[2]


%%
\begin{doc}{Règle de stabilité des configurations électroniques}
  Pour \important{augmenter leurs stabilités,} les atomes peuvent remplir leur couche externe en adoptant la configuration électronique du gaz noble avec le numéro atomique le plus proche.
  Ce principe se décompose en deux règles pour les 18 premiers éléments chimiques :
  
  \begin{importants}
    \pointCyan \important{Règle du duet :} les atomes de numéro atomique $Z < 6$
    tendent à adopter la configuration électronique
    \texteTrou[0.3]{de l'hélium : \important{1s$^2$}.}
    Ils ont \texteTrou[0.2]{2 (un duet)} électrons sur leur couche externe. 
    \bigskip
    
    \pointCyan \important{Règle de l'octet :} les atomes de numéro atomique $Z > 6$
    tendent à adopter la configuration électronique externe du gaz noble le plus proche
    \texteTrou(0){\important{ns$^2$ np$^6$}, avec $n = 2,3$.}
    Ils ont \texteTrou[0.2]{8 (un octet)} électrons sur leur couche externe.
  \end{importants}
\end{doc}


%%
\begin{doc}{Formation des ions}
  Pour adopter une configuration électronique plus stable, les atomes vont spontanément perdre ou gagner des électrons et ainsi former des ions avec une couche externe pleine.

  \hfill \textcolor{couleurSec}{\faLevelDown*}

  \begin{importants}
    \attention Dans la nature, on ne trouve jamais d'ions avec plus de 3 charges pour les éléments chimiques des 4 premières lignes du tableau périodiques, les éléments forment plutôt des molécules.
  \end{importants}

  \exemple L'aluminium a pour numéro atomique $Z = 13$, sa configuration électronique est 1s$^2$ 2s$^2$ 2p$^6$ 3s$^2$ 3p$^1$.
  Pour atteindre la configuration du néon avec 8 électrons sur la couche externe, il devra perdre 3 électrons pour former l'ion \chemfig{Al^{3+}}, avec 3 charges positives.
\end{doc}

\question{
  Le lithium \chemfig{Li} a pour numéro atomique $Z = 3$.
  Rappeler sa configuration électronique.
  Pour devenir stable, quelle règle doit-il respecter ? 
  Combien d'électrons doit-il perdre ou gagner pour la respecter ?
  Quel ion formera-t-il ?
}{
  Le lithium a 3 électrons, donc sa configuration électronique est $1s^2 2s^1$.
  Le gaz noble avec le numéro atomique le plus proche de lui est l'hélium, il va donc respecter la règle du duet et perdre 1 électrons. Il donc va former l'ion \chemfig{Li^+}.
}[5]


\question{
  Mêmes questions pour le soufre \isotope{S}{16}{} ($Z = 16$).
}{
  Le lithium a 16 électrons, donc sa configuration électronique est $1s^2 2s^2 2p^6 3s^2 3p^4$.
  Le gaz noble avec le numéro atomique le plus proche de lui est l'argon, il va donc respecter la règle de l'octet et gagner 2 électrons. Il donc va former l'ion \chemfig{S^{2-}}.
}[5]


\question{
  Par analogie avec le soufre \isotope{S}{16}{}, pouvez-vous répondre simplement aux mêmes questions pour l'oxygène \isotope{O}{8}{}, en utilisant leurs positions respectives dans le tableau périodique ?
}{
  Comme le soufre, il manque 2 électrons à l'oxygène pour respecter la règle de l'octet, il formera donc aussi l'ion \chemfig{O^{2-}}.
}[5]


\question{
  Comment répondre à ces questions en regardant simplement les colonnes du tableau périodique ?
}{
  Il suffit de compter le nombre de case d'écart qu'il y a entre l'élément chimique étudié et le gaz noble le plus proche.
}[5]